\ifdefined\gkver\else\def\gkver{student}\fi
\documentclass[\gkver]{gaokaozhenti}
\begin{document}

\chapter{2026年贵州}

\begin{enumerate}
\item
%% number: 1
%% paperName: 2026年贵州高考第 1 题 4 分
%% typeId: 1
%% type: 单选题
%% chapter: 直线运动
%% point: 匀变速直线运动
%% method: 无
%% score: 4
%% degree: 850
%% duplicateId: 0
%% body:
某同学乘坐列车时，在自己的座位上利用车厢内信息屏和手机秒表估算隧道长度。该同学进隧道时速度为 $60\Ums$，出隧道时速度为 $50\Ums$，总用时 $200\Us$。若列车在隧道中做匀减速直线运动，则该隧道长为 \xzanswer{C}

\fourchoices[answer=C]
{$5000\Um$}
{$8000\Um$}
{$11000\Um$}
{$14000\Um$}

%% answer: C

%% memo:
\memoanswer{
列车在隧道内做匀减速直线运动，匀变速直线运动的平均速度满足 $\overline{v}=\frac{v_{0}+v}{2}$。

位移等于平均速度乘以运动时间，即 $x=\overline{v}t$。

代入已知条件 $v_{0}=60\Ums$、$v=50\Ums$、$t=200\Us$。

计算得 $x=\frac{60+50}{2}\times200\Um=11000\Um$。

故选 C。
}
\item
%% number: 2
%% paperName: 2026年贵州高考第 2 题 4 分
%% typeId: 1
%% type: 单选题
%% chapter: 曲线运动
%% point: 万有引力定律
%% method: 无
%% score: 4
%% degree: 750
%% duplicateId: 0
%% body:
月球引潮力是引起海洋潮汐的主要原因，可等效为地表某点处质量为 $\Delta m$ 的海水所受月球引力减去地心处相同质量的物质所受月球引力。已知地球半径为 $R$，地心 $O$ 与月心 $O'$ 间距为 $d$，月球质量为 $M$，引力常量为 $G$，地表 $P$ 点背对月球，$P$、$O$ 和 $O'$ 在同一直线上，如图所示，则 $P$ 点处质量为 $\Delta m$ 的海水所受月球引潮力大小为 \xzanswer{D}
\onepicture[w=6.0cm]{figs/02.png}

\fourchoices[answer=D]
{$\left|\frac{GM\Delta m}{d}-\frac{GM\Delta m}{R}\right|$}
{$\left|\frac{GM\Delta m}{d^{2}}-\frac{GM\Delta m}{R^{2}}\right|$}
{$\left|\frac{GM\Delta m}{d+R}-\frac{GM\Delta m}{d}\right|$}
{$\left|\frac{GM\Delta m}{(d+R)^{2}}-\frac{GM\Delta m}{d^{2}}\right|$}

%% answer: D

%% memo:
\memoanswer{
根据题图可知 $P$ 点距月心的距离为 $d+R$，地心与月心间距为 $d$，月球质量为 $M$，根据万有引力公式结合月球引潮力的定义可得 $P$ 点处质量为 $\Delta m$ 的海水所受月球引潮力大小为 $F=\left|\frac{GM\Delta m}{(d+R)^{2}}-\frac{GM\Delta m}{d^{2}}\right|$。

故选 D。
}
\item
%% number: 3
%% paperName: 2026年贵州高考第 3 题 4 分
%% typeId: 1
%% type: 单选题
%% chapter: 运动和力
%% point: 动量和能量
%% method: 无
%% score: 4
%% degree: 550
%% duplicateId: 0
%% body:
如图，完全相同的均质小球 A、B 被不可伸长的细线悬挂，静止在同一竖直平面内，相互接触无挤压，悬挂点到球心的距离分别为 $L$ 和 $\frac{4}{5}L$，A 被拉至与竖直方向成 $60\Udu$ 的位置并由静止释放，随后与 B 发生弹性正碰。忽略空气阻力，B 的球心上升的最大高度为 \xzanswer{A}
\onepicture[w=3.6cm]{figs/03.png}

\fourchoices[answer=A]
{$\frac{1}{2}L$}
{$\frac{1}{3}L$}
{$\frac{1}{4}L$}
{$\frac{1}{5}L$}

%% answer: A

%% memo:
\memoanswer{
设 A、B 小球的质量均为 $m$，忽略空气阻力，则 A 从静止释放至与 B 发生碰撞前瞬间，由动能定理可得
$mgL(1-\cos60^{\circ})=\frac{1}{2}mv_{A}^{2}$

A 球与 B 球碰撞过程中，根据动量守恒定律和机械能守恒定律有 $mv_{A}=mv_{A}'+mv_{B}'$
$\frac{1}{2}mv_{A}^{2}=\frac{1}{2}mv_{A}'^{2}+\frac{1}{2}mv_{B}'^{2}$

可得 $v_{B}'=v_{A}$

由于 $\frac{1}{2}mv_{B}'^{2}<\frac{4}{5}mgL$

所以从碰撞后至下次碰撞前，B 做圆周运动，设 B 的球心上升的最大高度为 $h$，则对 B 从碰后至上升到最大高度的过程，由动能定理可得
$-mgh=0-\frac{1}{2}mv_{B}'^{2}$

解得 $h=\frac{1}{2}L$

故选 A。
}
\item
%% number: 4
%% paperName: 2026年贵州高考第 4 题 4 分
%% typeId: 1
%% type: 单选题
%% chapter: 电场
%% point: 带电粒子的偏转
%% method: 无
%% score: 4
%% degree: 450
%% duplicateId: 0
%% body:
如图，密闭真空中，有一竖直放置的金属靶和水平放置的两平行极板，极板与金属靶受光面垂直，板间存在竖直向上的匀强电场 $E$。用频率为 $\nu_{1}$ 和 $\nu_{2}$ 的光分别照射靶时，垂直靶面逸出最大初动能分别为 $E_{k1}$ 和 $E_{k2}$ 的光电子，经狭缝 $S_{1}$、$S_{2}$ 后进入电场，分别落到下极板 M、N 处。忽略极板边缘效应及电子间的相互作用，则 \xzanswer{A}
\onepicture[w=8.0cm]{figs/04.png}

\fourchoices[answer=A]
{$E_{k1}<E_{k2}$，$\nu_{1}<\nu_{2}$}
{$E_{k1}>E_{k2}$，$\nu_{1}>\nu_{2}$}
{$E_{k1}>E_{k2}$，$\nu_{1}<\nu_{2}$}
{$E_{k1}<E_{k2}$，$\nu_{1}>\nu_{2}$}

%% answer: A

%% memo:
\memoanswer{
光电子进入匀强电场后做类平抛运动，竖直方向：加速度 $a=\frac{qE}{m}$。

对所有电子，竖直方向加速度和位移均相同，根据 $y=\frac{1}{2}at^{2}$ 可知，运动时间相同。

水平方向做匀速运动，则落在 M 的电子初速度较小，即 $E_{k1}<E_{k2}$。

对光电效应过程，根据爱因斯坦光电效应方程 $E_{k}=h\nu-W_{0}$，可知 $\nu_{1}<\nu_{2}$。

故选 A。
}
\item
%% number: 5
%% paperName: 2026年贵州高考第 5 题 4 分
%% typeId: 1
%% type: 单选题
%% chapter: 电磁感应
%% point: 远距离输电
%% method: 无
%% score: 4
%% degree: 450
%% duplicateId: 0
%% body:
如图，某发电站输出电压 $U_{1}=13.8\UkV$ 的交流电，经两个变压器先后将电压升至 $U_{2}=220\UkV$ 和 $U_{3}=500\UkV$，再并入电网。输电线电阻 $R_{1}=12\UO$，$R_{2}=8\UO$。设所用变压器均为理想变压器，若发电站输出功率 $P_{1}=690\UkW$，则 $R_{1}$ 和 $R_{2}$ 上损失的总功率为 \xzanswer{B}
\onepicture[w=10.0cm]{figs/05.png}

\fourchoices[answer=B]
{$30.118\UkW$}
{$30.072\UkW$}
{$30.024\UkW$}
{$30\UkW$}

%% answer: B

%% memo:
\memoanswer{
发电站输出总功率 $P_{1}=U_{1}I_{1}$。

$I_{1}=\frac{P_{1}}{U_{1}}=\frac{690\times10^{3}\UW}{13.8\times10^{3}\UV}=50\UA$

因此 $R_{1}$ 所在回路电流：

$R_{1}$ 的功率损失：$\Delta P_{1}=I_{1}^{2}R_{1}=30000\UW=30\UkW$

理想变压器无损耗，第一个变压器副线圈输出功率等于原线圈输入功率，即 $P_{2}=P_{1}-\Delta P_{1}$

结合 $P_{2}=U_{2}I_{2}$，得 $R_{2}$ 所在回路电流：$I_{2}=3\UA$

$R_{2}$ 的功率损失：$\Delta P_{2}=I_{2}^{2}R_{2}=72\UW=0.072\UkW$

故选 B。
}
\item
%% number: 6
%% paperName: 2026年贵州高考第 6 题 4 分
%% typeId: 1
%% type: 单选题
%% chapter: 力物体的平衡
%% point: 三力平衡
%% method: 无
%% score: 4
%% degree: 550
%% duplicateId: 0
%% body:
如图，农民伯伯挑玉米时用双手分别抓住轻绳的 $O$、$O'$ 点处，某时刻绳 $OA$、$O'B$ 与竖直方向的夹角均为 $45\Udu$，手对绳的作用力分别垂直于 $OA$、$O'B$，所有作用力在同一竖直平面内，此时人和物均处于平衡状态。已知每筐玉米质量为 $m$，重力加速度大小为 $g$，忽略筐和扁担质量，则此时扁担对肩膀的作用力大小为 \xzanswer{B}
\onepicture[w=6.5cm]{figs/06.png}

\fourchoices[answer=B]
{$\frac{\sqrt{2}}{2}mg$}
{$mg$}
{$\sqrt{2}mg$}
{$2mg$}

%% answer: B

%% memo:
\memoanswer{
对单侧 O 点受力分析：O 点受三个力：向下的玉米重力 $mg$、$OA$ 绳的拉力 $T$、手垂直 $OA$ 的作用力 $F$，系统平衡。

根据平衡条件 $T=mg\cos45^{\circ}$

则单侧绳对扁担向下的分力为 $N_{1}=T\cos45^{\circ}$

扁担对肩膀的作用力大小为 $N=2N_{1}$

解得 $N=mg$

故选 B。
}
\item
%% number: 7
%% paperName: 2026年贵州高考第 7 题 4 分
%% typeId: 1
%% type: 单选题
%% chapter: 磁场
%% point: 带电粒子在复合场中的运动
%% method: 无
%% score: 4
%% degree: 450
%% duplicateId: 0
%% body:
如图，空间中存在垂直于纸面向里的匀强磁场及从 $O$ 点沿径向向外的电场，某处电场强度大小 $E=\frac{k}{r^{2}}$，$k$ 为常量，$r$ 为该处到 $O$ 的距离。一带负电粒子在纸面内沿逆时针方向做匀速圆周运动。当磁感应强度大小为 $B_{1}$ 时，粒子运动半径为 $r_{1}$，速率为 $v$，电势能为 $E_{p1}$；当磁感应强度大小为 $B_{2}$（$B_{2}>B_{1}$）时，粒子运动半径为 $r_{2}$，速率仍为 $v$，电势能为 $E_{p2}$。取无限远处的电势为零，不计粒子重力，则 \xzanswer{C}
\onepicture[w=4.3cm]{figs/07.png}

\fourchoices[answer=C]
{$r_{2}>r_{1}$，$E_{p2}>E_{p1}$}
{$r_{2}>r_{1}$，$E_{p2}<E_{p1}$}
{$r_{2}<r_{1}$，$E_{p2}<E_{p1}$}
{$r_{2}<r_{1}$，$E_{p2}>E_{p1}$}

%% answer: C

%% memo:
\memoanswer{
带负电粒子在纸面内沿逆时针方向做匀速圆周运动，则根据左手定则洛伦兹力背离圆心，电场力指向圆心，两力的合力提供向心力。

当磁感应强度大小为 $B_{1}$ 时，粒子运动半径为 $r_{1}$，速率为 $v$，此时有
$q\frac{k}{r_{1}^{2}}-qvB_{1}=\frac{mv^{2}}{r_{1}}$

当磁感应强度大小为 $B_{2}$（$B_{2}>B_{1}$）时，粒子运动半径为 $r_{2}$，速率仍为 $v$，此时有
$q\frac{k}{r_{2}^{2}}-qvB_{2}=\frac{mv^{2}}{r_{2}}$

联立可得 $q\frac{k}{r_{1}}-qvB_{1}r_{1}=q\frac{k}{r_{2}}-qvB_{2}r_{2}$

因为 $B_{2}>B_{1}$，所以 $r_{2}<r_{1}$

又沿电场方向电势降低，电子在电势低处电势能大，故 $E_{p2}<E_{p1}$。

故选 C。
}
\item
%% number: 8
%% paperName: 2026年贵州高考第 8 题 5 分
%% typeId: 2
%% type: 多选题
%% chapter: 气体性质
%% point: 热力学第一定律
%% method: 无
%% score: 5
%% degree: 650
%% duplicateId: 0
%% body:
在贵州凯里，人们常将小西红柿和红辣椒加工后放入如图所示的瓦罐中，罐口处倒扣一个钵并用水密封，发酵制作酸汤。发酵过程中罐内物质缓慢产生气体，压强足够大时气体以气泡形式溢出，间歇性放气。罐内气体可视为理想气体，设罐内气体温度和体积保持不变，则 \xzanswer{AD}
\onepicture[w=4.0cm]{figs/08.png}

\fourchoices[answer=AD]
{在放气过程中，溢出的气体对外界做正功}
{在放气过程中，溢出的气体对外界做负功}
{在刚放气完到下次即将放气的过程中，发酵产生气体，罐内气体内能减小}
{在刚放气完到下次即将放气的过程中，发酵产生气体，罐内气体内能增大}

%% answer: AD

%% memo:
\memoanswer{
AB．在放气过程中，气体体积增大，则溢出的气体对外界做正功，故 A 正确，B 错误；

CD．在刚放气完到下次即将放气的过程中，发酵产生气体，气体的质量增加，温度不变，罐内气体内能增大，故 C 错误，D 正确。

故选 AD。
}
\item
%% number: 9
%% paperName: 2026年贵州高考第 9 题 5 分
%% typeId: 2
%% type: 多选题
%% chapter: 光学
%% point: 光的干涉
%% method: 无
%% score: 5
%% degree: 450
%% duplicateId: 0
%% body:
光在折射率为 $n$ 的光纤中传播路程 $L$，等效于光在真空中传播路程 $nL$，$nL$ 称为光程。如图~\ref{2026贵州09a}，两根同种光纤在 A、B 处被熔接为耦合点，形成上下两支路，上下支路 A、B 之间的光纤长度相同。波长为 $\lambda$ 的光经过耦合点时的变化情况如图~\ref{2026贵州09b} 所示，任意支路的光经过耦合点，与输入口同侧的输出口的光未额外增加光程，异侧输出口的光额外增加 $\frac{\lambda}{4}$ 光程。现有波长为 $\lambda$ 的光从 1 口入射，则上下支路的光 \xzanswer{BC}
\twopicture[numA=a, numB=b, labelA=2026贵州09a, labelB=2026贵州09b, wA=7.0cm, wB=7.5cm]
{figs/09a.png}
{figs/09b.png}

\fourchoices[answer=BC]
{在 3 口干涉加强}
{在 3 口干涉削弱}
{在 4 口干涉加强}
{在 4 口干涉削弱}

%% answer: BC

%% memo:
\memoanswer{
现有波长为 $\lambda$ 的光从 1 口入射，经过耦合点 A 后，在上通道的光程为 $nL$，下通道的光程为 $nL+\frac{\lambda}{4}$，经过耦合点 B 后，上通道的光到 3 口光程为 $nL$，下通道的光到 3 口光程为 $nL+\frac{\lambda}{4}+\frac{\lambda}{4}$，两光的光程差为 $\frac{\lambda}{2}$，则在 3 口干涉削弱；经过耦合点 B 后，上通道的光到 4 口光程为 $nL+\frac{\lambda}{4}$，下通道的光到 4 口光程为 $nL+\frac{\lambda}{4}$，两光的光程差为零，则在 4 口干涉加强。

故选 BC。
}
\item
%% number: 10
%% paperName: 2026年贵州高考第 10 题 5 分
%% typeId: 2
%% type: 多选题
%% chapter: 电磁感应
%% point: 电磁感应与动量定理
%% method: 微元
%% score: 5
%% degree: 350
%% duplicateId: 0
%% body:
如图，水平绝缘桌面上固定有光滑金属导轨 ECADG，其中 EC 和 GD 的延长线交于 $\angle CAD$ 的角平分线上 O 点，M、N 点在 AO 的延长线上，M 为 CD 的中点，$\angle CAD=2\alpha$，$\angle EOG=2\beta$，$|OM|=d$，$|ON|=L$。导轨所在区域存在磁感应强度大小为 $B$、方向竖直向下的匀强磁场。质量为 $m$ 的导体棒在水平拉力作用下，从 A 点以速度 $v_{0}$ 向右匀速运动到 M 点，此时撤去拉力，运动到 N 点时速度为 $\frac{v_{0}}{2}$。已知导体棒单位长度电阻为 $r$，导体棒在运动过程中始终垂直于 AN 且与导轨接触良好，导轨电阻不计。则导体棒 \xzanswer{AC}
\onepicture[w=5.0cm]{figs/10.png}

\fourchoices[answer=AC]
{从 A 到 M 的过程中，电流为 $\frac{Bv_{0}}{r}$}
{从 M 到 N 的过程中，通过导体棒的电荷量为 $\frac{B(L^{2}-d^{2})}{2dr}$}
{速度 $v_{0}$ 满足关系式 $v_{0}=\frac{2B^{2}(L^{2}-d^{2})\tan\beta}{mr}$}
{从 A 到 M 的过程中，产生的焦耳热为 $\frac{B^{2}d^{2}v_{0}\tan\beta}{r\tan\alpha}$}

%% answer: AC

%% memo:
\memoanswer{
A．从 A 到 M 的过程中，导体棒向右匀速运动，设某时刻接入电路的导体棒长度为 $L'$，则有感应电动势 $E=BL'v_{0}$

导体棒接入电路的电阻为 $R'=L'r$

则电流为 $I=\frac{E}{R'}=\frac{Bv_{0}}{r}$，故 A 正确；

B．从 M 到 N 的过程中，结合 A 分析可知，某时刻导体棒速度为 $v$ 时，感应电流为 $I'=\frac{Bv}{r}$

由电流定义式可得，通过导体棒的电荷量为 $q=\overline{I}\Delta t=\sum I'\Delta t=\frac{B}{r}\sum v\Delta t$

又有 $\sum v\Delta t=L-d$

则有 $q=\frac{B(L-d)}{r}$，故 B 错误；

C．从 M 到 N 的过程中，令某时刻导体棒速度为 $v$ 时，导体棒的位移为 $x$，由几何关系可得，导体棒接入电路的长度为 $2(d+x)\tan\beta$，可知该时刻导体棒所受安培力为 $B\cdot\frac{Bv}{r}\cdot2(d+x)\tan\beta$

对该过程由动量定理有
$-\sum B\cdot\frac{Bv}{r}\cdot2(d+x)\tan\beta\cdot\Delta t=m\left(\frac{v_{0}}{2}-v_{0}\right)$

则有 $\frac{2B^{2}\tan\beta}{r}\sum(d+x)v\cdot\Delta t=\frac{mv_{0}}{2}$

解得 $v_{0}=\frac{4B^{2}\tan\beta}{mr}\sum(d+x)\cdot x$

作出该过程 $d+x$ 与 $x$ 关系图像，如图所示
\onepicture[w=4.5cm]{figs/10_an1.jpg}
则有 $\sum(d+x)\cdot x=\frac{1}{2}(d+L)(L-d)$

联立解得 $v_{0}=\frac{2B^{2}(L^{2}-d^{2})\tan\beta}{mr}$，故 C 正确；

D．从 A 到 M 的过程中，设某时刻导体棒位移大小为 $x'$，则导体棒接入电路的长度为 $2x'\tan\alpha$，该时刻导体棒所受安培力为 $B\cdot\frac{Bv_{0}}{r}\cdot2x'\tan\alpha$，该过程中安培力做功为
$W=-\sum B\cdot\frac{Bv_{0}}{r}\cdot2x'\tan\alpha\cdot x'=-\frac{2B^{2}v_{0}\tan\alpha}{r}\sum x'\cdot x'$

作出该过程 $x'$ 与 $x'$ 关系图像，如图所示
\onepicture[w=4.5cm]{figs/10_an2.jpg}
则有 $\sum x'\cdot x'=\frac{1}{2}d_{AM}^{2}=\frac{d^{2}\tan^{2}\beta}{2\tan^{2}\alpha}$

由功能关系有产生的焦耳热 $Q=|W|$

联立解得 $Q=\frac{B^{2}d^{2}v_{0}\tan^{2}\beta}{r\tan\alpha}$，故 D 错误。

故选 AC。
}
\item
%% number: 11
%% paperName: 2026年贵州高考第 11 题 6 分
%% typeId: 4
%% type: 实验题
%% chapter: 机械振动
%% point: 单摆测重力加速度
%% method: 无
%% score: 6
%% degree: 650
%% duplicateId: 0
%% body:
如图，某同学居家设计简易单摆测量重力加速度。主要器材有：细线、金属球、卷尺、手机、挂衣架。

请完成下列步骤：
\onepicture[h=4.0cm, align=right]{\includesvg{figs/11}}
\begin{enumerate}
	\item
	用卷尺测量金属球直径三次后得平均值 $d=1.82\Ucm$。
	\item
	用卷尺测量摆线长度三次后得平均值 $L$。将金属球从平衡位置拉至一个偏角小于 $5\Udu$ 的位置并由静止释放，用手机秒表测量金属球 $20$ 次全振动的时间 $t$。重复测量三次，得周期的平均值 $T$。
	\item
	改变摆线长度，重复步骤（2）。得 $L$、$t$（$t_{1}$、$t_{2}$、$t_{3}$）和 $T$ 的数据如下表：

	\begin{center}
	\begin{tabular}{|c|c|c|c|c|}
	\hline
	$L/\Ucm$ & $t_{1}/\Us$ & $t_{2}/\Us$ & $t_{3}/\Us$ & $T/\Us$ \\ \hline
	110.10 & 42.43 & 42.25 & 42.38 & 2.12 \\ \hline
	100.37 & 40.43 & 40.50 & 40.38 &  \\ \hline
	90.23 & 38.62 & 38.38 & 38.50 & 1.93 \\ \hline
	\end{tabular}
	\end{center}

	计算上表中 $L=100.37\Ucm$ 时的 $T=$ \tkanswer{2.02} $\Us$。（保留三位有效数字）
	\item
	根据重力加速度表达式 $g=$ \tkanswer{$\frac{2\pi^{2}(2L+d)}{T^{2}}$}（用 $\pi$、$T$、$L$ 和 $d$ 表示），计算实验值。
	\item
	查询当地重力加速度，发现 $g$ 的实验值偏小。为减小空气阻力引起的误差，应尽量选择质量 \tkanswer{大}（填“大”或“小”）、体积 \tkanswer{小}（填“大”或“小”）的摆球。
\end{enumerate}

%% answer:
\jdanswer{
\begin{enumerate}
	\item
	$d=1.82\Ucm$

	\item
	摆线长度 $L$，$20$ 次全振动的时间 $t$，周期 $T$

	\item
	$2.02$

	\item
	$\frac{2\pi^{2}(2L+d)}{T^{2}}$

	\item
	大，小

\end{enumerate}
}
%% memo:
\memoanswer{
（3）金属球做 $20$ 次全振动的时间为 $t_{1}+t_{2}+t_{3}$，则周期
$T=\frac{t_{1}+t_{2}+t_{3}}{20\times3}=\frac{t_{1}+t_{2}+t_{3}}{60}$

当 $L=100.37\Ucm$ 时，代入数据可得
$T=\frac{40.43+40.50+40.38}{60}\Us=2.02\Us$

（4）根据单摆的周期公式 $T=2\pi\sqrt{\frac{l}{g}}$，其中 $l=L+\frac{d}{2}$，则
$T=2\pi\sqrt{\frac{L+\frac{d}{2}}{g}}$

可得重力加速度表达式为 $g=\frac{2\pi^{2}(2L+d)}{T^{2}}$。

（5）为减小空气阻力引起的误差，应尽量选择质量大、体积小的摆球。
}
\item
%% number: 12
%% paperName: 2026年贵州高考第 12 题 9 分
%% typeId: 4
%% type: 实验题
%% chapter: 电路
%% point: 电路实验
%% method: 无
%% score: 9
%% degree: 450
%% duplicateId: 0
%% body:
在绝缘板上均匀涂敷一种导电材料，制作的电阻样品如图~\ref{2026贵州12a} 所示。所制三个电阻的涂层长度相同、宽度较小且相同、厚度不同（几微米），阻值均约为几十欧。某实验小组通过测量三个电阻阻值，比较其涂层厚度，所用器材有：
\onepicture[num=a, label=2026贵州12a, h=2.2cm, align=right]{figs/12a.png}

待测电阻 $R_{x1}$、$R_{x2}$ 和 $R_{x3}$；

电压表 V$_{1}$（量程 $0\sim3\UV$）；

电压表 V$_{2}$（量程 $0\sim5\UV$）；

电阻箱 $R_{0}$（阻值 $0\sim999.9\UO$）；

滑动变阻器 $R$（阻值 $0\sim10\UO$）；

直流电源 $E$（电动势 $4.5\UV$）；

开关；导线若干。

请完成下列步骤：
\begin{enumerate}
	\item
	实验小组设计了如图~\ref{2026贵州12b} 所示的电路图。根据图~\ref{2026贵州12b} 完成图~\ref{2026贵州12c} 中的实物图连线。

	\twopicture[numA=b, numB=c, labelA=2026贵州12b, labelB=2026贵州12c, h=5.5cm]
	{figs/12b.png}
	{figs/12c.png}
	\item
	接入 $R_{x1}$，闭合 S，调节 $R$ 和 $R_{0}$，$R_{0}$ 的示数如图~\ref{2026贵州12d} 所示，此时 $R_{0}$ 读数为 \tkanswer{36.2} $\UO$，电压表 V$_{1}$ 和 V$_{2}$ 读数分别为 $1.81\UV$ 和 $3.62\UV$，由此可知 $R_{x1}=$ \tkanswer{36.2} $\UO$。断开 S，改换待测电阻，重复上述步骤。分别测得 $R_{x2}=20.2\UO$，$R_{x3}=70.5\UO$。
	\onepicture[num=d, label=2026贵州12d, h=4.0cm, align=right]{\includesvg{figs/12d}}
	\item
	由三个电阻测量值可知 \tkanswer{$R_{x2}$}（填“$R_{x1}$”“$R_{x2}$”或“$R_{x3}$”）涂层最厚。
	\item
	实验完成后，该小组通过原理分析，发现此方法测电阻未考虑电压表 V$_{1}$ 内阻的影响。若已知电压表 V$_{1}$ 内阻为 $R_{V1}$，电压表 V$_{1}$ 和 V$_{2}$ 示数分别为 $U_{1}$ 和 $U_{2}$，电阻箱示数为 $R_{0}$，则 $R_{x}=$ \tkanswer{$\frac{U_{1}R_{0}R_{V1}}{R_{V1}U_{2}-(R_{0}+R_{V1})U_{1}}$}（用 $R_{V1}$、$R_{0}$、$U_{1}$ 和 $U_{2}$ 表示）。经分析发现，本实验中考虑电压表 V$_{1}$ 内阻后，不影响判断三个电阻涂层厚度关系。
\end{enumerate}

%% answer:
\jdanswer{
\begin{enumerate}
	\item
	如图

	\item
	$36.2\UO$，$36.2\UO$

	\item
	$R_{x2}$

	\item
	$\frac{U_{1}R_{0}R_{V1}}{R_{V1}U_{2}-(R_{0}+R_{V1})U_{1}}$

\end{enumerate}
}
%% memo:
\memoanswer{
（1）根据电路图可得实物连接图如图所示
\onepicture[w=5.5cm]{figs/12_an.png}

（2）由题图（d）可知 $R_{0}$ 的读数为 $36.2\UO$，结合题图（b）和串联电路分压规律 $\frac{R_{x}}{R_{x}+R_{0}}=\frac{U_{1}}{U_{2}}$

解得 $R_{x}=\frac{U_{1}}{U_{2}-U_{1}}R_{0}$

代入电压表的读数可得 $R_{x1}=\frac{1.81}{3.62-1.81}\times36.2\UO=36.2\UO$

（3）三个电阻导电材料相同、涂层长度相同、宽度较小且相同，由电阻定律 $R_{x}=\rho\frac{L}{S}$ 可知横截面积越大，即涂层越厚的电阻，阻值越小，由于 $R_{x2}$ 阻值最小，则 $R_{x2}$ 涂层最厚。

（4）若考虑电压表 V$_{1}$ 内阻，则根据串、并联电路规律和欧姆定律有
$\frac{U_{2}-U_{1}}{R_{0}}=\frac{U_{1}}{R_{x}}+\frac{U_{1}}{R_{V1}}$

解得 $R_{x}=\frac{U_{1}R_{0}R_{V1}}{R_{V1}U_{2}-(R_{0}+R_{V1})U_{1}}$
}
\item
%% number: 13
%% paperName: 2026年贵州高考第 13 题 9 分
%% typeId: 6
%% type: 计算题
%% chapter: 机械波
%% point: 波的图象
%% method: 无
%% score: 9
%% degree: 650
%% duplicateId: 0
%% body:
图~\ref{2026贵州13a} 是绳波在某时刻的照片。设该绳波为简谐横波，以绳中各质点平衡位置的连线为 $x$ 轴、质点振动方向为 $y$ 轴，建立如图~\ref{2026贵州13b} 所示的直角坐标系，图中实线、虚线分别表示 $t=0$ 和 $t=0.25\Us$ 时绳波的图像。若该绳波沿 $x$ 轴负方向传播，周期大于 $0.25\Us$。
\threepicture[numA=a, numB=b, numC=c, labelA=2026贵州13a, labelB=2026贵州13b, labelC=2026贵州13c, wA=4.3cm, wB=5.3cm, wC=4.6cm]
{figs/13a.png}
{figs/13b.png}
{figs/13c.png}
\begin{enumerate}
	\item
	求绳波的波长和振幅；
	\item
	求绳波的周期和传播速度；
	\item
	在图~\ref{2026贵州13c} 中，画出 $t=0.5\Us$ 时绳波在 $0\sim2.0\Um$ 内的图像。
\end{enumerate}

%% answer:
\jdanswer{
\begin{enumerate}
	\item
	$\lambda=2\Um$，$A=16\Ucm$

	\item
	$T=1\Us$，$v=2\Ums$

	\item
	如图

\end{enumerate}
}
%% memo:
\memoanswer{
（1）由题图（b）可知绳波的波长为 $\lambda=2\Um$，振幅 $A=16\Ucm$。

（2）由于绳波沿 $x$ 轴负方向传播，则由题图（b）可知绳波从 $t=0$ 到 $t=0.25\Us$，传播的时间为 $\Delta t=\frac{1}{4}T+nT$（$n=0,1,2,\cdots$）

由于 $T>0.25\Us$

则 $n$ 只能取 $0$，解得 $T=1\Us$

则波速 $v=\frac{\lambda}{T}=2\Ums$

（3）时间 $t=0.5\Us=\frac{T}{2}$，则此时的波形为 $t=0$ 时的图像沿 $x$ 轴负方向传播半个周期，$0\sim2.0\Um$ 内的图像为题图中实线关于 $x$ 轴对称的图像，如图所示
\onepicture[w=6.0cm]{figs/13_an.png}
}
\item
%% number: 14
%% paperName: 2026年贵州高考第 14 题 15 分
%% typeId: 1
%% type: 单选题
%% chapter: 电场
%% point: 带电粒子的直线运动
%% method: 无
%% score: 15
%% degree: 450
%% duplicateId: 0
%% body:
在竖直平面内，一带电荷量为 $q$（$q>0$）的小球在重力作用下从 P 点由静止开始下落，运动过程中始终受到与运动方向相反的空气阻力作用，其大小 $f$ 与速率 $v$ 满足 $f=kv$（$k$ 为常量）。小球第一次经过 P 点正下方的 M 点时达到最大速率 $v_{0}$，此时，施加竖直向上的恒定匀强电场，小球做变速运动。经过一段时间后，小球在 M 点正上方的 N 点再次达到最大速率 $v_{0}$，此后匀速上升。已知小球速率从第一次 $v_{0}$ 到再次达到 $v_{0}$ 的过程中，克服空气阻力做功为 $W_{f}$，重力加速度大小为 $g$。求：
\begin{enumerate}
	\item
	小球的质量和电场强度大小；
	\item
	小球的最大加速度大小；
	\item
	施加电场后，M、N 两点间的电势差。
\end{enumerate}

%% answer: E
\jdanswer{
\begin{enumerate}
	\item
	$m=\frac{kv_{0}}{g}$，$E=\frac{2kv_{0}}{q}$

	\item
	$a_{\max}=2g$

	\item
	$U_{MN}=\frac{2W_{f}}{q}$

\end{enumerate}
}
%% memo:
\memoanswer{
（1）小球第一次达到最大速率 $v_{0}$ 时，受力平衡，根据平衡条件有 $mg=kv_{0}$

解得小球的质量为 $m=\frac{kv_{0}}{g}$

小球再次到达最大速率 $v_{0}$ 时，受力平衡，小球到达 M 点上方 N 点，速度方向竖直向上，则空气阻力竖直向下，根据力的平衡条件 $mg+kv_{0}=qE$

解得电场强度大小为 $E=\frac{2kv_{0}}{q}$

（2）对小球从 P 点运动到 M 点的过程，由牛顿第二定律有 $mg-kv=ma$

可知小球做加速度减小的加速运动，从 M 点运动到最低点的过程，由牛顿第二定律有 $mg-qE+kv-mg=ma$

可知小球做加速度减小的减速运动，小球从最低点运动到 N 点的过程，由牛顿第二定律有 $qE-kv-mg=ma$

可知小球做加速度减小的加速运动，小球从静止释放时加速度为 $g$，小球运动到 M 点加恒定匀强电场时有 $qE+kv_{0}-mg=ma_{\max}$

解得 $a_{\max}=2g>g$

则小球的最大加速度大小为 $2g$。

（3）对小球从 M 点运动到 N 点的过程，由动能定理有 $qU_{MN}-mgh-W_{f}=0$

其中 $U_{MN}=Eh$

联立解得 M、N 两点间的电势差为 $U_{MN}=\frac{2W_{f}}{q}$。
}
\item
%% number: 15
%% paperName: 2026年贵州高考第 15 题 18 分
%% typeId: 6
%% type: 计算题
%% chapter: 曲线运动
%% point: 平抛运动
%% method: 无
%% score: 18
%% degree: 350
%% duplicateId: 0
%% body:
如图，机舱内装有质量均为 $m_{1}=5\Ukg$ 的甲、乙两个灭火弹模型的弹射型无人机，静止在 $\theta=30\Udu$ 的斜直轨道上，无人机空载时质量 $m_{0}=30\Ukg$。无人机在弹射系统作用下以 $v_{0}=15\Ums$ 的速度沿轨道离开，随后无人机依靠自身动力飞行，达到高度 $H=125\Um$ 时，开始以 $v_{1}=20\Ums$ 的速度沿水平方向做匀速直线运动，并进行投弹训练。设两弹所受空气阻力不计，落地点均在同一水平面上，取重力加速度 $g=10\Umsq$。
\onepicture[w=8.0cm, align=right]{figs/15.png}
\begin{enumerate}
	\item
	求载弹无人机在斜直轨道上运动过程中所受合力的冲量大小和方向；
	\item
	设水平飞行过程中，载弹无人机水平方向动力与质量满足 $F=\frac{1960}{m}$（国际单位制），所受空气阻力大小恒为 $f=49\UN$，方向与飞行方向相反，若两弹相对无人机无初速度先后被释放，时间间隔 $\Delta t=2\Us$，求两弹落地点之间的距离；
	\item
	设无人机水平飞行过程中，先相对无人机无初速度释放甲，当甲落地时沿水平方向发射乙，此时乙相对地面的速度大小为 $v$，若无人机的速度始终不变，求乙从发射到落地的过程中，两弹之间距离的最小值与 $v$ 取值的关系。
\end{enumerate}

%% answer:
\jdanswer{
\begin{enumerate}
	\item
	$600\UNs$，方向与水平方向夹角为 $30\Udu$

	\item
	$42.4\Um$

	\item
	$s_{\min}=\sqrt{125^{2}-\frac{(1250-v^{2})^{2}}{100}}$

\end{enumerate}
}
%% memo:
\memoanswer{
（1）载弹无人机在斜直轨道上运动过程中，由动量定理有 $I_{\text{合}}=(2m_{1}+m_{0})v_{0}$

代入数据解得 $I_{\text{合}}=600\UNs$

则载弹无人机在斜直轨道上运动过程中所受合力的冲量大小为 $600\UNs$，方向沿斜直轨道向上，即与水平方向夹角为 $30\Udu$。

（2）先释放一个灭火弹后载弹无人机水平方向的动力大小为 $F_{1}=\frac{1960}{m_{1}+m_{0}}\UN=56\UN$

水平方向由牛顿第二定律有 $F_{1}-f=(m_{1}+m_{0})a$

则释放第一个灭火弹后，第二个灭火弹与无人机一起做匀加速直线运动，所以后释放的灭火弹释放时的速度为 $v_{2}=v_{1}+a\Delta t$

从释放第一个灭火弹到释放第二个灭火弹的过程，由速度位移公式可知后释放的灭火弹的位移为 $x_{1}=\frac{v_{2}^{2}-v_{1}^{2}}{2a}$；由于两弹所受空气阻力不计，则两弹释放后均做平抛运动，有 $H=\frac{1}{2}gt_{0}^{2}$

则两弹落地点之间的距离为 $\Delta x=x_{1}+v_{2}t_{0}-v_{1}t_{0}=42.4\Um$

（3）相对无人机无初速度释放甲，且无人机的速度始终不变，则根据平抛运动规律可知，甲落地时，乙在甲的正上方 $H$ 处，设发射乙后乙的运动时间为 $t$，则此时两弹之间的距离为
$s=\sqrt{(vt)^{2}+\left(H-\frac{1}{2}gt^{2}\right)^{2}}=\sqrt{\frac{1}{4}g^{2}t^{4}+(v^{2}-gH)t^{2}+H^{2}}$

由二次函数知识可知，当 $t^{2}=\frac{2(gH-v^{2})}{g^{2}}$，两弹之间距离最小，即
$s_{\min}=\sqrt{H^{2}-\frac{(gH-v^{2})^{2}}{g^{2}}}$

由于 $t^{2}\ge0$，所以不考虑 $gH-v^{2}<0$

解得 $s_{\min}=\sqrt{125^{2}-\frac{(1250-v^{2})^{2}}{100}}$（国际单位制）
}
\end{enumerate}

\end{document}
